% Apresentação Beamer escrita para este estudo de caso.
% Compilação, na raiz: bash presentations/experimento/build.sh
\documentclass[aspectratio=169,10pt]{beamer}
\usepackage{fontspec}
\usepackage[brazilian]{babel}
\usepackage{graphicx,tikz,booktabs,array}
\usetikzlibrary{arrows.meta,positioning}
\setsansfont{texgyreheros-regular.otf}[BoldFont=texgyreheros-bold.otf,ItalicFont=texgyreheros-italic.otf]
\newfontfamily\displayfont{texgyrepagella-regular.otf}
\definecolor{Ink}{HTML}{173039}
\definecolor{Paper}{HTML}{F7F3EB}
\definecolor{Green}{HTML}{146B63}
\definecolor{Mist}{HTML}{BBD8CD}
\definecolor{Muted}{HTML}{657578}
\definecolor{Rust}{HTML}{AC553A}
\setbeamertemplate{navigation symbols}{}
\setbeamertemplate{headline}{}
\setbeamertemplate{footline}{}
\setbeameroption{hide notes}
\graphicspath{{presentations/experimento/assets/}}
\hypersetup{pdftitle={Uma dissertação pronta. Um pesquisador formado?},pdfauthor={Flavio Ribeiro},pdfsubject={Pesquisa acadêmica com agentes de IA},colorlinks=true,urlcolor=Green}
\setbeamercolor{normal text}{fg=Ink,bg=Paper}
\newenvironment{Canvas}{\begin{tikzpicture}[remember picture,overlay]\coordinate (O) at (current page.north west);}{\end{tikzpicture}}
% Texto em posições relativas à página, em centímetros. O texto continua editável.
\newcommand{\TextAt}[7]{\node[anchor=north west,inner sep=0pt,outer sep=0pt,text width=#3cm,align=left,text=#5,font={#6\fontsize{#4}{\fpeval{1.2*#4}}\selectfont}] at ([xshift=#1cm,yshift=-#2cm]O) {\hyphenpenalty=10000\exhyphenpenalty=10000 #7};}
\newcommand{\Body}[6]{\TextAt{#1}{#2}{#3}{#4}{#5}{\sffamily}{#6}}
\newcommand{\Bold}[6]{\TextAt{#1}{#2}{#3}{#4}{#5}{\sffamily\bfseries}{#6}}
\newcommand{\Display}[6]{\TextAt{#1}{#2}{#3}{#4}{#5}{\displayfont}{#6}}
\newcommand{\Header}[2]{\Bold{.9}{.48}{14}{8}{Green}{#1}\Bold{.9}{1.13}{14.2}{18}{Ink}{#2}}
\newcommand{\DarkHeader}[2]{\Bold{.9}{.48}{14}{8}{Mist}{#1}\Bold{.9}{1.13}{14.2}{18}{white}{#2}}
\newcommand{\Page}[1][Muted]{\Body{14.0}{8.52}{1.1}{8}{#1}{\insertframenumber/12}}
\newcommand{\Backup}{\Body{14.0}{8.52}{1.1}{8}{Muted}{Apoio}}
\newcommand{\CoverImage}[5]{\node[anchor=north west,inner sep=0pt] at ([xshift=#1cm,yshift=-#2cm]O) {\includegraphics[width=#3cm,height=#4cm,keepaspectratio]{#5}};}
\newcommand{\SpeakerNote}[1]{\note{\csname fala#1\endcsname}}
% Gerado por gerar_roteiro.py a partir de timing.json.
% Notas ocultas no PDF da audiência.

\expandafter\def\csname fala1\endcsname{Apresentar o experimento como uma investigação sobre a delegação da pesquisa e as evidências de formação. O tema científico serviu de terreno de teste. Não houve concessão de título, defesa ou aprovação institucional. A apresentação organiza uma discussão para a pós-graduação, sem pretender representar uma posição institucional.\par Fontes: Relatório do estudo de caso, abertura.}

\expandafter\def\csname fala2\endcsname{Texto e código, sozinhos, não atestam o domínio individual quando sua produção pode ser delegada. O experimento torna essa distinção concreta: o proponente declarou não dominar o tema e o agente entregou ambos. Isso não significa que documentos sejam inúteis para avaliar a pesquisa, nem que já tenham sido prova suficiente de aprendizagem. A pergunta é o que a pós-graduação passa a exigir da pessoa, além dos produtos. O caso não mediu aprendizagem.\par Fontes: Relatório, seção 3.}

\expandafter\def\csname fala3\endcsname{O proponente declarou não dominar o tema e não forneceu uma hipótese científica nova, equações ou código. Entregou um TG de 2003 e pediu uma continuação com potencial de publicação, implementação e validação. O texto do slide resume os pedidos, sem transcrevê-los. Houve direção humana, acompanhamento operacional e encaminhamento de pareceres. A busca de novidade começou por examinar a literatura: resultados do TG já tinham publicação posterior. Um antecedente recente motivou uma pergunta sobre a confiabilidade de simplificações. Originalidade era candidata, não certificada.\par Fontes: Relatório, seção 1 e apêndices A e B; TG original.}

\expandafter\def\csname fala4\endcsname{O diagrama resume a ordem predominante, com iterações entre as etapas. O usuário definiu a intenção. O Codex Astra formulou o recorte, organizou um plano de 13 marcos e coordenou três subagentes. Consultou antecedentes para evitar anunciar como novidade o que já existia. Executou quatro campanhas principais com 45.528 análises, que compartilham dados e não correspondem a igual número de simulações independentes. Os testes levaram a correções. Os dois pareceres do Gemini chegaram por intermédio do usuário. O artigo teve uma revisão após o encerramento do goal. O tempo humano não foi cronometrado. As três caixas escuras destacam as etapas conduzidas pelos agentes. Isso não significa ausência completa de participação humana nem permite atribuir uma porcentagem de autonomia.\par Fontes: Relatório, seções 1 e 2 e apêndices A a D.}

\expandafter\def\csname fala5\endcsname{Os quatro destaques são 167 páginas de dissertação, nove páginas de artigo revisado, 31 horas no contador do goal e US\$ 50 como rateio de assinatura estimado pelo proponente. O contador foi de 30h51min26s, com 54h56min42s no calendário, e não mede tempo de CPU. O artigo revisado é posterior ao goal. US\$ 50 não é uma fatura específica da execução. O equivalente em API Standard é cerca de US\$ 1.300, com os complementos científicos, conforme hipóteses e exclusões no apoio. As 26.069 linhas incluem comentários e vazios, e as 45.528 análises principais compartilham dados. Esses totais não medem qualidade nem equivalência a anos de trabalho humano. As capas aparecem no apoio.\par Fontes: PDFs da dissertação e do manuscrito; relatório, apêndices C a E.}

\expandafter\def\csname fala6\endcsname{O exemplo concreto é uma falha de calibração: para um parâmetro de ruído em 500 casos simulados, a faixa anunciada como 90\% continha a resposta conhecida em 313 casos, ou 62,6\%. Não se trata de intervalos fraudulentos nem de falha de todos os resultados do estudo. O agente registrou a limitação da aproximação. Em outros momentos, reconheceu que a comparação ampla continuava inconclusiva, respondeu a uma questão mais restrita e conteve alegações de originalidade. O slide reúne esses movimentos de revisão, sem atribuir todos à mesma causa. Não afirmar ausência geral de alucinações ou fabricação sem uma auditoria independente de todo o trabalho.\par Fontes: Relatório, seção 2 e apêndice B.1.}

\expandafter\def\csname fala7\endcsname{O Gemini elaborou dois pareceres, encaminhados pelo usuário. No primeiro ciclo, o Codex verificou críticas e acrescentou estudos pontuais, incluindo 300 análises em 50 casos. No segundo, revisou o manuscrito com os resultados já existentes. Aceitou melhorias de exposição e contexto, mas recusou alegações de originalidade absoluta e de falha inevitável em dados reais. As frases do slide são sínteses desses movimentos, não citações. Não houve revisão cega, certificação científica nem avaliação formal por pares humanos. Perguntar que trabalho de verificação permanece para orientador e revisores humanos.\par Fontes: Relatório, seção 2 e apêndice A.2; notas de revisão do artigo.}

\expandafter\def\csname fala8\endcsname{O título é uma provocação normativa para o CPG, não uma conclusão empírica de que a pesquisa incremental acabou. Defender que execução rotineira ou pequenas variações sem informação nova não bastam para justificar artigo ou título. Uma contribuição incremental pode revelar um regime novo, eliminar uma dúvida importante ou reproduzir criticamente um resultado. O debate é qual avanço de conhecimento exigir quando parte da execução pode ser delegada. O caso não prova colapso dos periódicos nem uma taxa de dez artigos por semana. Nenhuma fala privada ou pessoa é citada.\par Fontes: Relatório, seção 4.}

\expandafter\def\csname fala9\endcsname{Este caso foi computacional e não comparou sua produtividade à de um laboratório. Em simulações, controles com resposta conhecida permitem testar hipóteses. Medições acrescentam confronto com fenômenos e exigem instrumentação e calibração. O laboratório tem restrições materiais, mas também usa automação. Teoria e computação não perdem valor científico por serem automatizáveis. A discussão é como gerar evidência nova e interpretar seus limites, não proclamar superioridade de um tipo de pesquisa. A evidência do caso não demonstra desempenho em dados instrumentais reais. O destaque ao trabalho físico não pressupõe que toda operação de laboratório deva ser humana. Calibração e instrumentação também admitem automação.\par Fontes: Relatório, seção 4.}

\expandafter\def\csname fala10\endcsname{A expressão sobre execução virar commodity é uma provocação sobre a divisão do trabalho, sem afirmar que toda execução científica ficou fácil ou barata. Discutir como observar a capacidade de formular uma pergunta relevante, contestar uma saída convincente e rejeitar resultados sem evidência. Uma atividade possível é apresentar uma resposta do agente e pedir que o estudante explique quais sinais a enfraquecem e qual teste decidiria sua aceitação. Escrita e programação também são meios de aprendizagem e continuam exigindo domínio. A delegação abre oportunidades, mas o caso não mediu seu efeito sobre a formação.\par Fontes: Relatório, seções 3 e 4.}

\expandafter\def\csname fala11\endcsname{Proposta pedagógica, não intervenção avaliada neste caso. Antes de executar uma mudança, pedir uma previsão justificada. Depois, pedir que a pessoa explique o resultado, especialmente quando diverge da previsão. Na defesa, modificar uma hipótese ou condição, solicitar um argumento central e seus limites, com trechos sem assistência quando apropriado. Um teste ao vivo exige critérios justos e não deve reduzir domínio a rapidez de digitação ou memória. Registros ao longo do curso, qualificações e conversa oral podem complementar o documento final. Discutir como graduar profundidade e autonomia entre mestrado e doutorado.\par Fontes: Relatório, seções 3 e 5.}

\expandafter\def\csname fala12\endcsname{Reservar quase três minutos para abrir a discussão. A primeira questão trata da contribuição e dos patamares de originalidade e autonomia de mestrado e doutorado. A segunda trata de disciplinas, qualificações e defesas capazes de tornar compreensão observável. A terceira pergunta como ensinar a usar agentes sem impedir a construção de competências necessárias para julgá-los. O caso oferece artefatos e decisões auditáveis. Não estabelece uma política institucional nem demonstra aprendizagem. O relatório e os demais produtos estão no repositório, com fontes acessíveis no apoio.\par Fontes: Relatório, seção 5.}

\expandafter\def\csname fala13\endcsname{Números históricos da v1.0.0. O manuscrito de nove páginas é posterior. O total de código arquivado não é deduplicado. O tempo do goal não é CPU. As análises compartilham dados.\par Fontes: Relatório, apêndices C e D.}

\expandafter\def\csname fala14\endcsname{Aplicação dos preços oficiais Astra Standard de 13/09/2026 aos registros únicos por resposta. Até o goal são 6.510 respostas, com agente principal e três subagentes, totalizando US\$ 1.289,34. Os complementos científicos acrescentam US\$ 14,29. 97,56\% da entrada veio de cache. Gravações de cache zero no registro; cobrar toda entrada nova como gravação acrescentaria cerca de US\$ 53, mantida a reutilização. Fast dobra o preço. Não incluir o contador do goal nem somar raciocínio duas vezes. Exclui Gemini, computação local, ferramentas adicionais e estudo de caso. US\$ 50 é estimativa do usuário sobre rateio, não tarifa de API.\par Fontes: Relatório, apêndice E; estimativa\_custo\_api.json; developers.openai.com/api/docs/models/gpt-6-astra.}

\expandafter\def\csname fala15\endcsname{O repositório reúne PDFs e fontes, evidências e notas de revisão. O corpo do relatório discute decisões e formação; os apêndices guardam interações, métricas e método. Conversas privadas não foram transcritas. As imagens reproduzem as capas originais da dissertação e do manuscrito, com suas atribuições preservadas.\par Fontes: https://github.com/flavioluiz/TGdoWayne}

\begin{document}

% 01 — Capa. A questão precede os detalhes de execução.
\setbeamercolor{background canvas}{bg=Ink}
\begin{frame}[plain]
\begin{Canvas}
\Bold{.9}{.78}{14}{9}{Mist}{PESQUISA ACADÊMICA COM AGENTES DE IA}
\Display{.9}{2.40}{14.2}{30}{white}{Uma dissertação pronta.}
\Display{.9}{3.65}{14.2}{36}{Mist}{Um pesquisador\\formado?}
\Body{.94}{6.91}{14}{12}{white}{Um experimento de pesquisa delegada e revisão entre modelos}
\Bold{.94}{8.05}{5}{11}{white}{Flavio Ribeiro}
\Body{8.3}{8.10}{6.8}{9}{Mist}{Pós-graduação no ITA / 13 setembro 2026}
\end{Canvas}
\SpeakerNote{1}
\end{frame}

% 02 — A pergunta sobre formação já aparece na abertura.
\setbeamercolor{background canvas}{bg=Paper}
\begin{frame}[plain]
\begin{Canvas}
\Header{A PERGUNTA}{O documento e a formação}
\Display{.9}{2.70}{14.2}{30}{Ink}{Um bom documento\\basta para demonstrar\\domínio individual?}
\Body{.94}{6.65}{14}{14}{Green}{O caso permite discutir a qualidade da entrega\\e a capacidade de explicar e sustentar suas escolhas.}
\Page
\end{Canvas}
\SpeakerNote{2}
\end{frame}

% 03 — A capa é evidência documental, com atribuição original.
\begin{frame}[plain]
\begin{Canvas}
\Header{O EXPERIMENTO}{O desafio inicial}
\CoverImage{.95}{2.32}{3.8}{5.42}{tg-01.png}
\Bold{5.7}{2.62}{9.3}{12}{Green}{Um trabalho de graduação de 2003}
\Display{5.7}{3.75}{9.3}{21}{Ink}{Propor uma continuação.\\Implementar e validar.\\Escrever a dissertação.}
\Body{5.7}{6.96}{9.1}{12}{Muted}{O proponente não dominava o tema\\e não forneceu uma hipótese nova.}
\Page
\end{Canvas}
\SpeakerNote{3}
\end{frame}

% 04 — Diagrama nativo. A ordem resume um processo com iterações.
\begin{frame}[plain]
\begin{Canvas}
\Header{O EXPERIMENTO}{O percurso da pesquisa delegada}
\tikzset{etapa/.style={draw=Green,line width=.75pt,fill=Paper,text=Ink,minimum width=3.8cm,minimum height=1.1cm,inner sep=5pt,font=\sffamily\bfseries\fontsize{10.5}{13}\selectfont},fluxo/.style={Green,line width=1.2pt,-{Stealth[length=5pt]}}}
\node[etapa] (a) at ([xshift=2.8cm,yshift=-3.3cm]O) {Instrução ampla};
\node[etapa] (b) at ([xshift=8cm,yshift=-3.3cm]O) {Busca e recorte};
\node[etapa] (c) at ([xshift=13.2cm,yshift=-3.3cm]O) {Execução};
\node[etapa] (d) at ([xshift=13.2cm,yshift=-5.6cm]O) {Testes e correções};
\node[etapa] (e) at ([xshift=8cm,yshift=-5.6cm]O) {Revisão pelo Gemini};
\node[etapa] (f) at ([xshift=2.8cm,yshift=-5.6cm]O) {Dissertação e artigo};
\draw[fluxo] (a)--(b);\draw[fluxo] (b)--(c);\draw[fluxo] (c)--(d);\draw[fluxo] (d)--(e);\draw[fluxo] (e)--(f);
\Bold{.94}{7.22}{14.1}{12}{Green}{Codex Astra, com três subagentes.\\Pareceres encaminhados pelo usuário.}
\Page
\end{Canvas}
\SpeakerNote{4}
\end{frame}

% 05 — Volume, tempo e custo têm definições detalhadas nos slides de apoio.
\begin{frame}[plain]
\begin{Canvas}
\Header{O EXPERIMENTO}{O que o agente entregou}
\CoverImage{.95}{2.32}{3.5}{4.55}{dissertacao-001.png}
\CoverImage{5.15}{2.32}{3.5}{4.55}{artigo-1.png}
\Bold{.95}{7.08}{3.6}{13}{Ink}{Dissertação}
\Body{.95}{7.67}{4.1}{9.5}{Muted}{167 páginas\\47 referências}
\Bold{5.15}{7.08}{3.6}{13}{Ink}{Manuscrito}
\Body{5.15}{7.67}{4.1}{9.5}{Muted}{9 páginas\\Após a revisão}
\Display{10.1}{2.40}{4.9}{42}{Green}{31 h}
\Body{10.16}{3.92}{4.8}{13}{Ink}{no contador do goal}
\Body{10.16}{4.70}{4.8}{12}{Ink}{55 h no calendário}
\Bold{10.16}{6.22}{4.8}{18}{Green}{\textasciitilde US\$ 1.300}
\Body{10.16}{7.11}{4.8}{11}{Muted}{API Standard estimada}
\Page
\end{Canvas}
\SpeakerNote{5}
\end{frame}

% 06 — Dois episódios de correção, sem estatística especializada.
\begin{frame}[plain]
\begin{Canvas}
\Header{O EXPERIMENTO}{Os testes mudaram as conclusões}
\Bold{.94}{2.52}{6.4}{9}{Muted}{O QUE OS TESTES MOSTRARAM}
\Bold{8.60}{2.52}{6.5}{9}{Muted}{COMO O TRABALHO MUDOU}
\Display{.94}{3.47}{6.45}{20}{Rust}{Certeza excessiva\\sobre o resultado}
\Body{8.60}{3.51}{6.45}{14}{Ink}{O texto explicitou\\a limitação do método.}
\Display{.94}{5.68}{6.45}{20}{Rust}{A pergunta ampla\\ficou inconclusiva}
\Body{8.60}{5.74}{6.45}{14}{Ink}{Um teste adicional\\respondeu a uma\\pergunta mais restrita.}
\Page
\end{Canvas}
\SpeakerNote{6}
\end{frame}

% 07 — Exemplos de aceitação e recusa de sugestões de outro modelo.
\setbeamercolor{background canvas}{bg=Ink}
\begin{frame}[plain]
\begin{Canvas}
\DarkHeader{O EXPERIMENTO}{O Gemini criticou. O Codex conferiu.}
\Bold{.94}{2.56}{6.5}{10}{Mist}{CRÍTICAS INCORPORADAS}
\Bold{8.6}{2.56}{6.5}{10}{Mist}{CONCLUSÕES RECUSADAS}
\Display{.94}{3.55}{6.5}{17}{white}{Novos testes.\\Mais contexto\\para o artigo.}
\Display{8.6}{3.55}{6.5}{17}{white}{Originalidade absoluta.\\Falha inevitável\\em dados reais.}
\Bold{.94}{7.17}{14.1}{13}{Mist}{O que orientadores e revisores precisam verificar?}
\Page[Mist]
\end{Canvas}
\SpeakerNote{7}
\end{frame}

% 08 — Cenário para debate, sem previsão de colapso editorial.
\setbeamercolor{background canvas}{bg=Paper}
\begin{frame}[plain]
\begin{Canvas}
\Header{IMPLICAÇÕES PARA A PESQUISA}{Qual contribuição merece atenção?}
\Display{.94}{2.70}{14.1}{30}{Ink}{O que passamos a saber\\que antes não sabíamos?}
\Body{.94}{5.56}{14.1}{17}{Green}{Uma extensão pequena pode resolver\\uma incerteza importante.}
\Body{.94}{7.63}{14.1}{11.5}{Muted}{Produzir mais manuscritos também exige mais capacidade de avaliação.}
\Page
\end{Canvas}
\SpeakerNote{8}
\end{frame}

% 09 — Computação e laboratório: diferenças de evidência, sem hierarquia.
\begin{frame}[plain]
\begin{Canvas}
\Header{IMPLICAÇÕES PARA A PESQUISA}{Simular e medir exigem evidências distintas}
\Display{.94}{2.74}{6.45}{25}{Green}{Simulação}
\Display{8.6}{2.74}{6.45}{25}{Green}{Medição}
\Body{.94}{4.15}{6.45}{15}{Ink}{Testar hipóteses\\com resposta conhecida.}
\Body{8.6}{4.15}{6.45}{15}{Ink}{Confrontar o modelo\\com o fenômeno.}
\Body{.94}{6.98}{14.1}{14}{Muted}{Em ambos os casos, alguém precisa justificar\\a ligação entre a evidência e a conclusão.}
\Page
\end{Canvas}
\SpeakerNote{9}
\end{frame}

% 10 — Oportunidade de executar mais e necessidade de aprender a julgar.
\begin{frame}[plain]
\begin{Canvas}
\Header{FORMAÇÃO E AVALIAÇÃO}{O que precisamos aprender a fazer?}
\Bold{.94}{2.53}{14.1}{10}{Muted}{O AGENTE AMPLIA A EXECUÇÃO}
\Display{.94}{3.58}{14.1}{26}{Ink}{A formação precisa desenvolver\\o julgamento sobre o trabalho.}
\Body{.94}{6.31}{14.1}{17}{Green}{Por que esta pergunta importa?\\O que faria você desconfiar do resultado?}
\Page
\end{Canvas}
\SpeakerNote{10}
\end{frame}

% 11 — Proposta pedagógica, não resultado testado pelo experimento.
\begin{frame}[plain]
\begin{Canvas}
\Header{FORMAÇÃO E AVALIAÇÃO}{A avaliação pode acompanhar as decisões}
\Bold{.94}{2.80}{4.1}{10}{Green}{ANTES DO TESTE}
\Bold{6.08}{2.80}{4.3}{10}{Green}{DURANTE A PESQUISA}
\Bold{11.24}{2.80}{3.9}{10}{Green}{NA DEFESA}
\Display{.94}{3.80}{4.1}{25}{Ink}{Prever}
\Display{6.08}{3.80}{4.3}{24}{Ink}{Interpretar}
\Display{11.24}{3.80}{3.9}{24}{Ink}{Sustentar}
\Body{.94}{5.14}{4.1}{14}{Ink}{O que mudará?\\Por quê?}
\Body{6.08}{5.14}{4.3}{14}{Ink}{Por que houve\\surpresa?}
\Body{11.24}{5.14}{3.9}{13.5}{Ink}{O argumento\\resiste à\\alteração?}
\Bold{.94}{7.56}{14.1}{12}{Green}{Domínio demonstrado ao longo do curso, além do PDF.}
\Page
\end{Canvas}
\SpeakerNote{11}
\end{frame}

% 12 — Encerramento aberto à discussão.
\setbeamercolor{background canvas}{bg=Ink}
\begin{frame}[plain]
\begin{Canvas}
\DarkHeader{DISCUSSÃO}{Três questões para a pós-graduação}
\Body{.94}{2.75}{1}{17}{Mist}{01}
\Display{2.58}{2.64}{12.5}{20}{white}{Que originalidade exigiremos\\no mestrado e no doutorado?}
\Body{.94}{4.68}{1}{17}{Mist}{02}
\Display{2.58}{4.57}{12.5}{20}{white}{Como disciplinas e qualificações\\podem demonstrar compreensão?}
\Body{.94}{6.61}{1}{17}{Mist}{03}
\Display{2.58}{6.50}{12.5}{20}{white}{Como ensinar a delegar pesquisa\\e a julgar o que o agente entrega?}
\Page[Mist]
\end{Canvas}
\SpeakerNote{12}
\end{frame}

% Apoio — A partir daqui, fora dos 25 minutos sugeridos.
\appendix
\setbeamercolor{background canvas}{bg=Paper}
\begin{frame}[plain]
\begin{Canvas}
\Header{EVIDÊNCIAS E FONTES}{Definições das métricas}
\node[anchor=north west,inner sep=0pt,text=Ink] at ([xshift=.94cm,yshift=-2.7cm]O) {\sffamily\fontsize{11}{14}\selectfont
\renewcommand{\arraystretch}{1.8}\begin{tabular}{@{}p{3.5cm}p{10.25cm}@{}}
\toprule
\textbf{Medida} & \textbf{Leitura correta}\\\midrule
26.069 linhas & Módulos, scripts e testes, com comentários e vazios\\
45.528 análises & Quatro campanhas, com dados compartilhados\\
54 h 56 min & Calendário da criação ao encerramento do goal\\
30 h 51 min & Contador do goal, sem medir CPU ou tempo de formação\\
\bottomrule
\end{tabular}};
\Backup
\end{Canvas}
\SpeakerNote{13}
\end{frame}

\begin{frame}[plain]
\begin{Canvas}
\Header{EVIDÊNCIAS E FONTES}{Hipóteses da estimativa de custo}
\node[anchor=north west,inner sep=0pt,text=Ink] at ([xshift=.94cm,yshift=-2.65cm]O) {\sffamily\fontsize{11}{14}\selectfont
\renewcommand{\arraystretch}{1.4}\begin{tabular}{@{}p{6.5cm}p{4.15cm}r@{}}
\toprule
\textbf{Categoria até o goal} & \textbf{Milhões de tokens} & \textbf{US\$}\\\midrule
Entrada nova & 20,58 & 205,84\\
Entrada em cache & 822,05 & 822,05\\
Saída, com raciocínio & 5,23 & 261,46\\\bottomrule
\end{tabular}};
\Bold{.94}{6.25}{14.1}{16}{Green}{\textasciitilde US\$ 1.300 Standard / \textasciitilde US\$ 2.600 Fast}
\Body{.94}{7.04}{14.1}{11.5}{Muted}{\textasciitilde US\$ 50 no rateio da assinatura informado pelo proponente.}
\Body{.94}{7.73}{14.1}{10}{Muted}{API com complementos científicos. Estimativas, não faturas.}
\Backup
\end{Canvas}
\SpeakerNote{14}
\end{frame}

\begin{frame}[plain]
\begin{Canvas}
\Header{EVIDÊNCIAS E FONTES}{Documentos para examinar o caso}
\Display{.94}{2.76}{14.1}{27}{Ink}{Um caso aberto ao exame.}
\Body{.94}{4.27}{14.1}{16}{Ink}{Relatório com evidências e apêndices\\Dissertação e manuscrito revisado\\Código, resultados e registros de revisão}
\Bold{.94}{7.15}{14.1}{14}{Green}{\href{https://github.com/flavioluiz/TGdoWayne}{Abrir o repositório e os materiais}}
\Backup
\end{Canvas}
\SpeakerNote{15}
\end{frame}
\end{document}
